\documentclass[numbers=noenddot, a4paper, 12pt, headsepline, footsepline]{scrreprt}

% Deutsche Umlaute
\usepackage[utf8]{inputenc}

% deutsche Silbentrennung
\usepackage[ngerman]{babel}

%anführungszeichen
\usepackage[autostyle=true,german=quotes]{csquotes}

% Einbinden von Bildern
\usepackage{graphicx}

% Mathematische Pakete
\usepackage{amsmath}
\usepackage{amssymb}
\usepackage{amsfonts}
\usepackage{amsthm}
\usepackage{latexsym}
\usepackage{morefloats}
\usepackage{geometry} 


\usepackage{standalone}

% better looking tables
\usepackage{ifthen}
\usepackage{booktabs}
\usepackage{multirow}


% ausgabe von quelltext
\usepackage{textcomp}
\usepackage{listings}

%einbinden von pdf seiten
\usepackage{pdfpages}

\lstset{
	%	backgroundcolor=\color{lbcolor},
	tabsize=4,    
	%	rulecolor=,
	language=[GNU]C++,
	basicstyle=\scriptsize,
	upquote=true,
	%aboveskip={1.5\baselineskip},
	columns=fixed,
	showstringspaces=false,
	extendedchars=false,
	breaklines=true,
	%prebreak = \raisebox{0ex}[0ex][0ex]{\ensuremath{\hookleftarrow}},
	%frame=single,
	numbers=none,
	showtabs=false,
	showspaces=false,
	showstringspaces=false,
	identifierstyle=\ttfamily,
	language=C++,
	basicstyle=\ttfamily,
	keywordstyle=\color{blue}\ttfamily,
	stringstyle=\color{red}\ttfamily,
	commentstyle=\color[rgb]{0,.7,0}\ttfamily,
	morecomment=[l][\color{magenta}]{\#},
	%\lstdefinestyle{C++}{language=C++,style=numbers}’,
}
\lstset{
	basicstyle=\renewcommand{\baselinestretch}{.9}\ttfamily,
	xleftmargin=.5cm,
	language=C++,
	numbers=left,
	stringstyle=\color{black}\ttfamily,
	keywordstyle=\color{black}\ttfamily,
	numberstyle=\tiny}


% coolere referenzen
\usepackage[german]{fancyref}

%literatur mittels biber
\usepackage[backend=biber,
			style=alphabetic
			%citestyle=alphabetic 
			]{biblatex}

%\ExecuteBibliographyOptions{
%	sorting=nyt, %Sortierung Autor, Titel, Jahr
%	bibwarn=true, %Probleme mit den Daten, die Backend betreffen anzeigen
%}
\addbibresource{literatur.bib}

%querformat seiten
\usepackage{lscape}

%überschrift des Kapitel auf jeder Seite
\pagestyle{headings}

%lange tabellen mit seitenumbruch
\usepackage{longtable}

%abkürzungsverzeichnis
\usepackage{acronym}

%um links im pdf zu haben mit denen man hin und her springen kann
\usepackage{hyperref}

%%text farbig hinterlegen
%\usepackage{color}

%%um text durch, unter, über... streichen zu können
%\usepackage{ulem}

%um in tabellen zellen diagonal teilen zu können
%\usepackage{diagbox}


\begin{document}

%einbinden titelblatt
\title{Smart Car} 
\subtitle{Hauptseminar Projektstudium}
\author{\\
	\begin{tabular}{|c|c|c|}
		\hline 
		Knorr Thomas 		 & xyzk 	& HSP1 \\ 
		\hline 
		Lackner Andreas 	 & xyzk 	& HSP2 \\ 
		\hline 
		Schleinkofer Michael & xyzk 	& HSP2 \\ 
		\hline 
		Welker Franz 		 & 3119754  & HSP2 \\ 
		\hline 
		Wiche Fabian 		 & xyzk		& HSP1 \\ 
		\hline 
	\end{tabular} 
	}	

\date{\today}
%\clearpairofpagestyles

\makeatletter
\begin{titlepage}
	
	\begin{center}
		\includegraphics[width=10cm]{./Images/OTHLogo.eps}\\
		\vspace{4cm}
		{\huge\bfseries\@title\unskip\strut\par}\paragraph{}
		{\Large\bfseries\@subtitle\unskip\strut\par}\paragraph{}
		An der\\
		Ostbayerischen Technischen Hochschule Regensburg\\
		Fakultät Informatik/Mathematik
		\paragraph{}
		Betreut durch Prof. Dr. Klaus Volbert\\
		\vspace{\fill}
		Vorgelegt von: \vspace{0,3cm}
		\@author\\
		\vspace{0,3cm}
		Datum: \@date
	\end{center}
\end{titlepage}
\makeatother



%arabische zahlen für seiten
\pagenumbering{arabic}

% Inhaltsverzeichnis anzeigen
%\parskip am schluss anpassen, so dass ein zweiseitiges inhaltsverzeichnis gut umgebrochen wird
{\parskip=+5mm \tableofcontents}

\chapter*{Einleitung}
 
Kurze Einleitung, warum haben wir uns das Thema ausgesucht (nicht lang oder so)...
Gibt z.B. Pace aber braucht man immer Handy, wollten es anders machen...











\chapter{Use Cases}

Was haben wir uns vorgestellt, dass wir am Ende erreichen wollen. 
Z.B. einfaches wie Strecken im Backend anzeigen mit infos wie länge, V,...
Im Dongle aufnehmen können
In der App live Daten sehen können

  

 
 
 
 
 
 
 
 % Warum haben wir es getan
\chapter{Implementierung}

Wie wurde es umgesetzt, wie waren die Ideen (also hier z.b. auf die Gesamte Architektur eingehen, spezifische der einzelnen Bereiche in den unter Kapiteln)

\section{App}

Was kann die App, und besonderes erklären

  

 
 
 
 
 
 
 
 
\section{Backend}

Was kann das Backend, und besonderes erklären
  

 
 
 
 
 
 
 
 
\section{Dongle}

Was kann der Dongle, und besonderes erklären

 
 
 
 
 
 
 
 

  

 
 
 
 
 
 
 
 % Wie haben wir es getan + aufgetretene Probleme
\chapter{Tests}

Was haben wir getestet, fahren aufgenommen und verifiziert....

  

 
 
 
 
 
 
 
 % Was haben wir getestet
\chapter{Ausblick}

Was kann man jetzt mit unserem Zeug als Grundlage anstellen, was haben wir uns überlegt aber konnten es nicht mehr umsetzen.

  

 
 
 
 
 
 
 
 % Was kann man jetzt mit unserer Grundlage weiterhin machen






\newpage
\listoffigures
\lstlistoflistings
\listoftables
% Muss von Hand sortiert werden!!

%anlegen:
% \acro{}[]{}
% \acro{}{}

%verwenden
%\ac{x} dann wird x als verlinkung genommen  
%\acp{x} wird auch x als verlinkung genommen, aber mit einem s am ende für mehrzahl

%\newpage
\section*{Abkürzungsverzeichnis} %Abkürzungsverzeichnis
\markboth{Abkürzungsverzeichnis}{}
\begin{acronym}[ABCDEFGHIJK] 			% längste Abkürzung steht in eckigen Klammern
	\acro{GPS}{Global Positioning System}
	\acro{MEMS}{Microelectromechanical System}
	\acro{NMEA}{National Marine Electronics Association}
  	\acro{API}{Application Programming Interface}
	\acro{OBD}{On Board Diagnistic}
	\acro{PID}{Parameter Identifikator}
	\acro{SBAS}{Satellite-based augmentation systems}
	\acro{UART}{Universal asynchronous receiver-transmitter}
	\acro{VIN}{Vehicle Identification Number}
	\acro{BLE}{Bluetooth Low Energy}
	\acro{MVID}{Mapped Vehilce Identification Number}
	\acro{RAM}{Random Access Memory}
	\acro{LINQ}{Language Integrated Query}
\end{acronym}

\printbibliography

\newpage
\appendix










\end{document}